\documentclass[10pt,a4paper]{amsart}
\usepackage[margin=19mm]{geometry}
\usepackage{amsmath,amssymb,mathtools}
\usepackage[scr=rsfs,cal=euler]{mathalfa}
\usepackage{xfrac}
\usepackage[unicode,hidelinks,bookmarks=false]{hyperref}
\newcommand{\ie}{i.e.\ }
\newcommand{\cf}{cf.\ }
\newcommand{\resp}{respectively\ }
\newcommand{\Subsection}{\subsection}
\providecommand{\nobreakdash}{\nobreak\hbox{-}}
\setlength{\emergencystretch}{2em}
\multlinegap0pt
\newtheorem{lem}{Lemma}
\title{On the Cauchy problem for the Langevin-type fractional equation}
\author{Yusuf Fayziev, Shakhnoza Jumaeva}
\date{Source: arXiv:2505.08252v1 (CC BY 4.0)}
\begin{document}
\maketitle

\noindent\emph{Source and licence.} On the Cauchy problem for the Langevin-type fractional equation. Version arXiv:2505.08252v1 (CC BY 4.0), distributed by arXiv under the Creative Commons Attribution 4.0 International licence, http://creativecommons.org/licenses/by/4.0/, as recorded in the arXiv licence metadata for this exact version and shown on the abstract page https://arxiv.org/abs/2505.08252v1. Full licence notice: \texttt{licenses/arxiv-2505.08252-CC-BY-4.0.txt}.\par\smallskip

\noindent\textit{Editorial scope.} Counted unit: the article's lem Lemma6 (source0.tex lines 243--248). The article's complete original proof of it is delivered with it (source0.tex lines 249--264, one \texttt{proof} environment of 16 lines). No mathematics is written, repaired or re-derived: the statement and the proof are the article's own lines, in the article's own notation, and no retained line is substituted or re-flowed.  Retained uncounted as the source-local context the proof needs: lines 180--186.  Nothing was omitted from any retained span, so the package carries no omission record.  No retained line contains a citation, so the wrapper carries no bibliography and no bibliography entry.  No retained line contains a figure environment, an image-inclusion command or drawing code, so no image is delivered and the package declares no asset dependency.  Editorial formatting only, in addition: an AMS \texttt{amsart} preamble that restates the article's own declarations and macros used by the retained lines, namely the environments \texttt{lem} on separate counters, together with the standard packages those lines need.  The retained environments print in this document as Lemma 1 in order of appearance, whereas the article prints the same items with the numbering of its own full document.  The complete native source of the article as distributed in the arXiv e-print for this exact version is bundled unchanged as \texttt{sources/178024/source0.tex}.
\begin{lem}\label{Lemma2}
  The following relation holds:
$$
|t^{\alpha-1}E_{\alpha,\mu}(-\lambda t^\alpha)|\leq C_\varepsilon \lambda^{\varepsilon-1}t^{\varepsilon\alpha-1} ,
$$
where $\lambda$ is a positive number and $0<\varepsilon<1$.
\end{lem}

\begin{lem}\label{Lemma6}
Let $0<\varepsilon<1$ be any fixed number, and $f(t) \in C([0,T];D(A^\varepsilon))$. Then the following estimate holds:
\begin{equation}\label{epsiloneq}
  \sum_{k=1}^\infty \left| \lambda_k\int_0^t (t-\eta)^{\alpha-1} E_{\alpha,\mu}(-\lambda_k (t-\eta)^\alpha)f_k(\eta) d\eta\right|^2\leq C_\varepsilon \max_{t\in[0,T]}||f||_\varepsilon^2.
\end{equation}
\end{lem}

\begin{proof}
 By using Lemma \ref{Lemma2} for any fixed number $0<\varepsilon<1$, we take
$$
\sum_{k=1}^n \lambda_k^2 \left|\int_0^t (t-\eta)^{\alpha-1}E_{\alpha,\mu}(-\lambda_k (t-\eta)^\alpha)f_k(\eta) d\eta\right|^2\leq
$$
$$\leq C^1_\varepsilon\sum_{k=1}^n  \left[\int_0^t (t-\eta)^{\varepsilon\alpha-1}\lambda_k^\varepsilon |f_k(\eta)| d\eta\right]^2.$$
Using the generalized Minkowski inequality, we have
$$C^1_\varepsilon\sum_{k=1}^n  \left[\int_0^t (t-\eta)^{\varepsilon\alpha-1}\lambda_k^\varepsilon |f_k(\eta)| d\eta\right]^2\leq C^1_\varepsilon\left(\int_0^t (t-\eta)^{\varepsilon \alpha-1}\bigg(\sum_{k=1}^n \lambda_k^{2\varepsilon} |f_k(\eta)|^2\bigg)^{\frac{1}{2}} d\eta\right)^2
$$
$$
\leq C^1_\varepsilon T^{\varepsilon\alpha}\max_{t\in[0,T]} ||f||_\varepsilon^2=C_\varepsilon \max_{t \in[0;T]}||f||_\varepsilon^2.
$$
Taking the limit as $n \rightarrow \infty$, we obtain the estimate \eqref{epsiloneq}.

Lemma \ref{Lemma6} has been proved.
\end{proof}

\end{document}
